\documentclass[11pt,letterpaper,oneside]{report}
\usepackage[T1]{fontenc}
\usepackage[utf8]{inputenc}
\usepackage[english]{babel}
\usepackage{lmodern}
\usepackage{microtype}
\usepackage[margin=0.86in,headheight=16pt]{geometry}
\usepackage[table,dvipsnames]{xcolor}
\usepackage{graphicx}
\usepackage{booktabs,tabularx,longtable,array}
\usepackage{enumitem}
\usepackage{titlesec}
\usepackage{fancyhdr}
\usepackage{hyperref}
\usepackage{xurl}
\Urlmuskip=0mu plus 1mu\relax
\usepackage{tcolorbox}
\usepackage{listings}
\usepackage{tikz}
\usetikzlibrary{arrows.meta,positioning,shapes.geometric,fit,calc}
\usepackage{needspace}
\usepackage{amssymb}
\definecolor{ink}{HTML}{14253A}
\definecolor{navy}{HTML}{17385D}
\definecolor{blue}{HTML}{20739A}
\definecolor{pale}{HTML}{ECF4F8}
\definecolor{green}{HTML}{1D6C5C}
\definecolor{softgreen}{HTML}{EEF8F4}
\definecolor{border}{HTML}{D4DFE7}
\definecolor{grey}{HTML}{506174}
\hypersetup{colorlinks=true,linkcolor=navy,urlcolor=blue,pdfauthor={Micro-Frontend Engineering Documentation Suite}}
\titleformat{\chapter}[display]{\normalfont\sffamily\bfseries\color{ink}}{\Large\color{blue}\MakeUppercase{\chaptername} \thechapter}{8pt}{\Huge}
\titlespacing*{\chapter}{0pt}{-10pt}{24pt}
\titleformat{\section}{\Large\sffamily\bfseries\color{navy}}{\thesection}{0.7em}{}
\titleformat{\subsection}{\large\sffamily\bfseries\color{navy}}{\thesubsection}{0.6em}{}
\setlist[itemize]{itemsep=3pt,topsep=5pt,leftmargin=1.6em}
\setlist[enumerate]{itemsep=4pt,topsep=5pt,leftmargin=1.8em}
\pagestyle{fancy}\fancyhf{}
\fancyhead[L]{\footnotesize\sffamily MICRO-FRONTEND ENGINEERING SUITE}
\fancyhead[R]{\footnotesize\sffamily\docshort}
\fancyfoot[L]{\footnotesize\color{grey}Engineering edition | October 2026}
\fancyfoot[R]{\footnotesize\thepage}
\renewcommand{\headrulewidth}{0.35pt}
\renewcommand{\arraystretch}{1.2}
\setlength{\emergencystretch}{2em}
\lstdefinestyle{eng}{basicstyle=\ttfamily\footnotesize,keywordstyle=\color{navy}\bfseries,commentstyle=\color{grey},stringstyle=\color{green},breaklines=true,breakatwhitespace=false,columns=fullflexible,keepspaces=true,showstringspaces=false,frame=single,rulecolor=\color{border},backgroundcolor=\color{pale},xleftmargin=5pt,xrightmargin=5pt,aboveskip=10pt,belowskip=10pt,captionpos=b}
\lstdefinelanguage{JavaScript}{morekeywords={const,let,var,require,return,async,await,import,export,default,function,try,catch,new,true,false,undefined,module},sensitive=true,morecomment=[l]//,morestring=[b]',morestring=[b]"}
\lstset{style=eng}
\newtcolorbox{sourcebox}{colback=pale,colframe=border,arc=2pt,boxrule=0.6pt,left=8pt,right=8pt,top=7pt,bottom=7pt}
\newtcolorbox{engineerbox}{colback=softgreen,colframe=green!45,arc=2pt,boxrule=0.6pt,left=8pt,right=8pt,top=7pt,bottom=7pt}
\newcommand{\based}[1]{\begin{sourcebox}\textbf{BOOK-GROUNDED CONCEPT}\\#1\end{sourcebox}}
\newcommand{\added}[1]{\begin{engineerbox}\textbf{ENGINEERING EXTENSION}\\#1\end{engineerbox}}
\newcommand{\term}[1]{\textbf{\color{navy}#1}}
\newcommand{\criterion}[1]{\textbf{Acceptance:} #1}
\newcommand{\book}{\emph{Building Micro-Frontends}, 2nd ed.}
\newcommand{\refs}{\chapter*{References and Provenance}\addcontentsline{toc}{chapter}{References and Provenance}
\begin{enumerate}[leftmargin=*]
\item \textbf{Primary source:} Luca Mezzalira, \emph{Building Micro-Frontends: Distributed Systems for the Frontend}, 2nd ed., O'Reilly Media, published November 2025 (copyright 2026), ISBN 978-1-098-17078-3. Chapter and printed-page references are included throughout. This suite paraphrases and synthesizes concepts; it does not reproduce the book.
\item \textbf{Engineering extension:} Webpack Module Federation reference: \url{https://webpack.js.org/plugins/module-federation-plugin}.
\item \textbf{Engineering extension:} Next.js multi-zones: \url{https://nextjs.org/docs/app/guides/multi-zones}. Details are version-sensitive and must be checked against the deployed framework version.
\item \textbf{Engineering extension:} OWASP Application Security Verification Standard 5.0: \url{https://owasp.org/www-project-application-security-verification-standard/}.
\item \textbf{Engineering extension:} GitHub Actions secure use reference: \url{https://docs.github.com/en/actions/reference/security/secure-use}.
\end{enumerate}}

\newcommand{\docshort}{PRACTICAL LABS}
\begin{document}
\begin{titlepage}\thispagestyle{empty}
\vspace*{1.1in}{\Huge\sffamily\bfseries\color{navy} Hands-On Engineering Laboratories\par}
\vspace{18pt}{\Large\sffamily\color{blue} Twelve progressive labs, runnable scaffolds, acceptance tests, and capstone\par}
\vspace{27pt}\noindent\color{blue}\rule{0.75\textwidth}{1.4pt}
\vspace{23pt}\par{\large\sffamily A companion to Luca Mezzalira's \emph{Building Micro-Frontends}, Second Edition}\par
\vfill\noindent\textbf{Document series:} Micro-Frontend Software Architecture and Engineering\\
\textbf{Version:} 1.0 \quad \textbf{Prepared:} October 2026\\
\textbf{Production format:} pdfLaTeX-compatible standalone source\\
\textbf{Scope:} Original synthesis plus explicitly identified engineering extensions
\end{titlepage}
\pagenumbering{roman}\tableofcontents\clearpage\pagenumbering{arabic}
\renewcommand{\chaptername}{LAB}\chapter*{Lab Environment and Assessment Rules}\addcontentsline{toc}{chapter}{Lab Environment and Assessment Rules}
\based{The labs reinforce concepts from Chapters 2--12 and organizational practices from Chapter 13. These are newly authored exercises, not the book's downloadable code samples or preexisting labs.}
\section*{Suggested environment}
Node.js LTS supported by the selected framework, npm/pnpm with lockfile, TypeScript, Webpack 5 or an equivalent runtime-composition experiment, an HTTP reverse proxy, a browser automation framework, Git, and a CI service. For security exercises, use a non-production isolated project and synthetic identities only. Select an actual supported Node/framework version and document it in a lockfile; no single package version is prescribed by the book.
\section*{Evidence convention}
Each submission includes a README with setup and cleanup instructions, a tested revision ID, diagrams/ADRs, passing/failing test outputs, screenshot or trace of relevant runtime behavior, and a risk note. Deliver scripts rather than screenshots alone wherever feasible.
\section*{Rubric per lab}
\begin{tabularx}{\linewidth}{p{.25\linewidth}X}\toprule Dimension & Weight and grading question \\\midrule
Correctness & 35: Does the implementation exhibit the intended behavior?\\
Reproducibility & 20: Can another engineer run and verify it?\\
Boundary fidelity & 20: Is the business and deployment contract explicit?\\
Security and resilience & 15: Are negative paths and failures tested?\\
Trade-off analysis & 10: Are alternatives and constraints documented?\\\bottomrule\end{tabularx}

\chapter{Domain mapping workshop}
\section{Prerequisites and objective}
\textbf{Prerequisites:} Basic HTML/JS; product-domain description.\par
\textbf{Outcome:} Map business capabilities, user flows, source-code ownership, and deployment boundaries.
\section{Procedure}
Collect routes and features; group by business outcomes; draw context map; annotate owner and API dependencies.
\section{Submission artifacts}
Capability map, candidate bounded contexts, ADR-001.
\section{Acceptance tests}
\criterion{Every proposed micro-frontend represents a testable business capability; two rejected alternatives documented.} Also run at least one deliberately failing test and document its expected failure mode.
\section{Extension challenge}
Re-run the same experiment with a different architecture or with an injected network/performance constraint; report the changed outcome and whether the initial recommendation still holds.

\chapter{Baseline a frontend monolith}
\section{Prerequisites and objective}
\textbf{Prerequisites:} Local Node runtime and browser test runner.\par
\textbf{Outcome:} Establish a quantitative pre-decomposition baseline.
\section{Procedure}
Build a small SPA; record dependency graph, cold load, JS transferred, CI duration, keyboard access, and release procedure.
\section{Submission artifacts}
Baseline report with scripts and browser traces.
\section{Acceptance tests}
\criterion{Measurements reproducible on a documented machine/network profile.} Also run at least one deliberately failing test and document its expected failure mode.
\section{Extension challenge}
Re-run the same experiment with a different architecture or with an injected network/performance constraint; report the changed outcome and whether the initial recommendation still holds.

\chapter{Route-owned vertical decomposition}
\section{Prerequisites and objective}
\textbf{Prerequisites:} Labs 1--2.\par
\textbf{Outcome:} Implement domain routing without runtime same-view composition.
\section{Procedure}
Implement shell/router; split catalog and account routes; direct-link/refresh; test 404 and unavailable zone behavior.
\section{Submission artifacts}
Two independently built applications, route map, routing tests.
\section{Acceptance tests}
\criterion{Each route has one owner; unrelated app can deploy independently.} Also run at least one deliberately failing test and document its expected failure mode.
\section{Extension challenge}
Re-run the same experiment with a different architecture or with an injected network/performance constraint; report the changed outcome and whether the initial recommendation still holds.

\chapter{Runtime Module Federation}
\section{Prerequisites and objective}
\textbf{Prerequisites:} Labs 1--3; Webpack 5.\par
\textbf{Outcome:} Compose one catalog remote with the shell at runtime.
\section{Procedure}
Configure remote exposes and host remotes; version dependency policy; load remote and fallback; test missing entry and chunk.
\section{Submission artifacts}
Remote and shell build logs, module contract, fallback screenshot.
\section{Acceptance tests}
\criterion{Remote unavailable does not crash global navigation; contract test passes.} Also run at least one deliberately failing test and document its expected failure mode.
\section{Extension challenge}
Re-run the same experiment with a different architecture or with an injected network/performance constraint; report the changed outcome and whether the initial recommendation still holds.

\chapter{Web component compatibility}
\section{Prerequisites and objective}
\textbf{Prerequisites:} DOM events and custom elements.\par
\textbf{Outcome:} Test an alternative independent integration surface.
\section{Procedure}
Implement a custom element with a documented input, emit a typed signal, verify CSS scoping and keyboard interaction.
\section{Submission artifacts}
Custom element, event schema, a11y test.
\section{Acceptance tests}
\criterion{Two different hosts can mount with no undocumented state dependency.} Also run at least one deliberately failing test and document its expected failure mode.
\section{Extension challenge}
Re-run the same experiment with a different architecture or with an injected network/performance constraint; report the changed outcome and whether the initial recommendation still holds.

\chapter{SSR zone routing and cache}
\section{Prerequisites and objective}
\textbf{Prerequisites:} Basic HTTP, proxy, web SSR.\par
\textbf{Outcome:} Compare full-document route composition with client federation.
\section{Procedure}
Route /account and /catalog to independently deployed origins; measure navigation; test auth and cache policy.
\section{Submission artifacts}
Proxy config, response headers, timing comparison.
\section{Acceptance tests}
\criterion{Private response not served to wrong session; deep links stable.} Also run at least one deliberately failing test and document its expected failure mode.
\section{Extension challenge}
Re-run the same experiment with a different architecture or with an injected network/performance constraint; report the changed outcome and whether the initial recommendation still holds.

\chapter{BFF and authorization}
\section{Prerequisites and objective}
\textbf{Prerequisites:} Backend basics, REST API and sessions.\par
\textbf{Outcome:} Create a trusted API boundary for a UI domain.
\section{Procedure}
Implement identity-aware BFF; schema validation; upstream timeout; resource-level authorization; security tests.
\section{Submission artifacts}
OpenAPI schema, negative tests, sequence diagram.
\section{Acceptance tests}
\criterion{Unauthorized record access forbidden server-side; no token in URL or custom event.} Also run at least one deliberately failing test and document its expected failure mode.
\section{Extension challenge}
Re-run the same experiment with a different architecture or with an injected network/performance constraint; report the changed outcome and whether the initial recommendation still holds.

\chapter{Consumer-driven contract tests}
\section{Prerequisites and objective}
\textbf{Prerequisites:} Labs 3--7.\par
\textbf{Outcome:} Catch incompatible remote and API changes before rollout.
\section{Procedure}
Create shared event schema fixtures; test current and previous contract version; introduce breaking change deliberately.
\section{Submission artifacts}
CI contract checks and documented version policy.
\section{Acceptance tests}
\criterion{Breaking event/API change fails prior to deployment.} Also run at least one deliberately failing test and document its expected failure mode.
\section{Extension challenge}
Re-run the same experiment with a different architecture or with an injected network/performance constraint; report the changed outcome and whether the initial recommendation still holds.

\chapter{Security controls and dependency audit}
\section{Prerequisites and objective}
\textbf{Prerequisites:} GitHub Actions or equivalent.\par
\textbf{Outcome:} Integrate secure build, SAST/CodeQL where supported, secret and dependency checks.
\section{Procedure}
Set least-privilege workflow; lock dependencies; review action provenance; test intentional insecure fixture without shipping it.
\section{Submission artifacts}
Threat model, pipeline outputs, exception register.
\section{Acceptance tests}
\criterion{No secrets exposed to untrusted PR; finding triage is traceable.} Also run at least one deliberately failing test and document its expected failure mode.
\section{Extension challenge}
Re-run the same experiment with a different architecture or with an injected network/performance constraint; report the changed outcome and whether the initial recommendation still holds.

\chapter{Frontend discovery and rollout}
\section{Prerequisites and objective}
\textbf{Prerequisites:} Labs 3--9.\par
\textbf{Outcome:} Promote immutable releases without rebuilding or modifying domain source.
\section{Procedure}
Implement controlled manifest; publish two artifact versions; 10% canary then rollback; trace selected release.
\section{Submission artifacts}
Manifest history, staged rollout, rollback evidence.
\section{Acceptance tests}
\criterion{Rollback restores prior release; unknown remote origin is denied.} Also run at least one deliberately failing test and document its expected failure mode.
\section{Extension challenge}
Re-run the same experiment with a different architecture or with an injected network/performance constraint; report the changed outcome and whether the initial recommendation still holds.

\chapter{Observe and recover a partial outage}
\section{Prerequisites and objective}
\textbf{Prerequisites:} Metrics/traces and failure injection.\par
\textbf{Outcome:} Instrument browser-to-API user journey and practice rollback.
\section{Procedure}
Tag telemetry with route/domain/release; inject remote unavailable/latency; inspect alert; invoke runbook.
\section{Submission artifacts}
Timeline, dashboard captures, postmortem.
\section{Acceptance tests}
\criterion{Fault attributed to domain/release; recovery confirmed by user-flow test.} Also run at least one deliberately failing test and document its expected failure mode.
\section{Extension challenge}
Re-run the same experiment with a different architecture or with an injected network/performance constraint; report the changed outcome and whether the initial recommendation still holds.

\chapter{Architecture review and migration capstone}
\section{Prerequisites and objective}
\textbf{Prerequisites:} Labs 1--11.\par
\textbf{Outcome:} Defend a full adoption choice and migration plan with evidence.
\section{Procedure}
Evaluate modular SPA vs route zones vs federation; record four decisions; design a 3-wave migration; demonstrate successful and failed rollout.
\section{Submission artifacts}
Architecture packet, ADRs, scorecard, pilot and retirement criteria.
\section{Acceptance tests}
\criterion{All four decisions addressed; risk, UX, security, performance, governance and rollback justified.} Also run at least one deliberately failing test and document its expected failure mode.
\section{Extension challenge}
Re-run the same experiment with a different architecture or with an injected network/performance constraint; report the changed outcome and whether the initial recommendation still holds.

\chapter*{Starter Contract and Test Fixtures}\addcontentsline{toc}{chapter}{Starter Contract and Test Fixtures}
\section*{Minimal domain descriptor}
\begin{lstlisting}[caption={Illustrative descriptor}]
export interface DomainDescriptor {
  id: 'catalog' | 'account' | 'checkout';
  routePrefix: string;
  buildRevision: string;
  contractMajor: number;
  remoteUrl?: string;
}
\end{lstlisting}
\section*{Policy validation pseudocode}
\begin{lstlisting}[caption={Pseudocode; implement with actual schema validation}]
function acceptDomain(candidate, policy): boolean {
  return policy.allowedDomains.includes(candidate.id)
    && policy.routeOwners[candidate.routePrefix] === candidate.id
    && candidate.contractMajor === policy.requiredContractMajor
    && (candidate.remoteUrl === undefined
        || policy.allowedOrigins.includes(new URL(candidate.remoteUrl).origin));
}
\end{lstlisting}
\noindent Origin allowlisting is only one layer. Require an authorized promotion service and verified immutable artifact identity before a remote is trusted. A string matching \texttt{https://} does not establish code provenance.
\section*{Failure injection matrix}
\begin{tabularx}{\linewidth}{p{.34\linewidth}X}\toprule Injection & Expected evidence \\\midrule
Remote host DNS failure & Shell remains navigable; domain fallback emitted\\
Manifest points to old remote & Previous contract supported or safe fail\\
User session expires mid-flow & Protected API rejects; login journey recoverable\\
BFF returns 503 & Domain-specific error state; correlation ID logged\\
CSS token version skew & Visual regression detects unintended appearance\\
Canary experiences higher error rate & Release paused or rolled back by documented procedure\\\bottomrule\end{tabularx}
\chapter*{Capstone Deliverable Checklist}\addcontentsline{toc}{chapter}{Capstone Deliverable Checklist}
\begin{itemize}[label=$\square$]\item 3 architecture alternatives and scored trade-off sheet.
\item 4 core decisions captured as accepted ADRs.
\item At least 2 independently built domains and 1 working shell/route entry.
\item Schema-tested cross-domain event and API contract.
\item Authentication and authorization negative-test evidence.
\item Accessible cross-domain keyboard journey and baseline web-vital report.
\item CI verification, artifact digest, discovery record, and canary rollback.
\item Trace/log evidence identifying faulted domain and version.
\item Migration plan, ownership matrix, and explicit non-goals.\end{itemize}

\refs
\end{document}
